\chapter{L17 -- Recap su decomposizione di Kalman, esercizi su raggiungibilità/osservabilità e forma minima}
\label{chap:L17}

\section{Obiettivo della lezione}

Questa lezione fa da raccordo tra i risultati già visti su:
\begin{itemize}
    \item sottospazio di raggiungibilità \(\mathcal R\),
    \item sottospazio di inosservabilità \(\bar{\mathcal O}\),
    \item forma standard di raggiungibilità,
    \item forma standard di osservabilità,
    \item decomposizione di Kalman,
\end{itemize}
e una serie di esercizi nei quali questi concetti vengono usati in modo operativo.

L'idea di fondo è sempre la stessa:
\begin{quote}
\emph{la dinamica interna del sistema può essere scomposta in parti che l'ingresso riesce o non riesce ad eccitare, e in parti che l'uscita riesce o non riesce a vedere.}
\end{quote}

In particolare, questa lezione insiste su tre aspetti:
\begin{enumerate}
    \item come si costruisce concretamente l'intersezione \(\mathcal R\cap\bar{\mathcal O}\);
    \item come si riconoscono i modi raggiungibili e osservabili in esercizi concreti;
    \item perché la funzione di trasferimento ``vede'' solo la parte minima, cioè quella sia raggiungibile sia osservabile.
\end{enumerate}

\bigskip

\section{Richiamo: costruzione dell'intersezione \(\mathcal R\cap\bar{\mathcal O}\)}

Riprendiamo un esempio già discusso nella lezione precedente, nel quale erano note due basi:

\[
T_R=
\begin{bmatrix}
2 & -1\\
3 & -3\\
1 & -1
\end{bmatrix},
\qquad
T_{\bar O}=
\begin{bmatrix}
2\\
0\\
0
\end{bmatrix}.
\]

Le colonne di \(T_R\) generano il sottospazio di raggiungibilità \(\mathcal R\), mentre le colonne di \(T_{\bar O}\) generano il sottospazio di inosservabilità \(\bar{\mathcal O}\).

Vogliamo trovare
\[
\mathcal R\cap\bar{\mathcal O}.
\]

\subsection{Metodo generale}

Un vettore \(v\) appartiene all'intersezione se e solo se può essere scritto sia come combinazione lineare delle colonne di \(T_R\), sia come combinazione lineare delle colonne di \(T_{\bar O}\). Dunque dobbiamo imporre

\[
T_R y = T_{\bar O} z,
\]
dove in questo caso
\[
y=
\begin{bmatrix}
y_1\\y_2
\end{bmatrix},
\qquad
z=
\begin{bmatrix}
z_1
\end{bmatrix}.
\]

Equivalentemente:

\[
\begin{bmatrix}
T_R & -T_{\bar O}
\end{bmatrix}
\begin{bmatrix}
y\\z
\end{bmatrix}
=0.
\]

Sostituendo i dati:

\[
\begin{bmatrix}
2 & -1 & -2\\
3 & -3 & 0\\
1 & -1 & 0
\end{bmatrix}
\begin{bmatrix}
y_1\\y_2\\z_1
\end{bmatrix}
=0.
\]

Il sistema lineare associato è

\[
\begin{cases}
2y_1-y_2-2z_1=0,\\
3y_1-3y_2=0,\\
y_1-y_2=0.
\end{cases}
\]

Le ultime due equazioni danno subito
\[
y_1=y_2.
\]
Sostituendo nella prima si ottiene
\[
2y_1-y_1-2z_1=0
\quad\Longrightarrow\quad
z_1=\frac{y_1}{2}.
\]

Scegliendo, ad esempio, \(y_1=y_2=1\), si ha
\[
z_1=\frac12.
\]

Quindi un vettore dell'intersezione è

\[
v=T_R
\begin{bmatrix}
1\\1
\end{bmatrix}
=
\begin{bmatrix}
2 & -1\\
3 & -3\\
1 & -1
\end{bmatrix}
\begin{bmatrix}
1\\1
\end{bmatrix}
=
\begin{bmatrix}
1\\0\\0
\end{bmatrix}.
\]

D'altra parte

\[
T_{\bar O}\frac12
=
\begin{bmatrix}
2\\0\\0
\end{bmatrix}\frac12
=
\begin{bmatrix}
1\\0\\0
\end{bmatrix},
\]
quindi è tutto coerente.

\subsection{Conclusione}

Abbiamo trovato

\[
\mathcal R\cap\bar{\mathcal O}
=
\operatorname{span}\left\{
\begin{bmatrix}
1\\0\\0
\end{bmatrix}
\right\}.
\]

Questa procedura è importante perché nella decomposizione di Kalman bisogna prima individuare proprio questa intersezione, cioè la parte del sistema che è:
\begin{itemize}
    \item raggiungibile,
    \item ma non osservabile.
\end{itemize}

\bigskip

\section{Esercizio 1: prova del 18/12/2023, sistema SISO a tempo discreto}

Consideriamo il sistema a tempo discreto

\[
\Sigma:
\begin{cases}
x(k+1)=Ax(k)+Bu(k),\\
y(k)=Cx(k)+Du(k),
\end{cases}
\]
con

\[
A=
\begin{bmatrix}
0 & 1 & 1 & 0\\
0 & 1 & 0 & 1\\
1 & 0 & 0 & 1\\
1 & 0 & 1 & 0
\end{bmatrix},
\qquad
B=
\begin{bmatrix}
1\\
1\\
0\\
0
\end{bmatrix},
\qquad
C=
\begin{bmatrix}
0 & 1 & 0 & 0
\end{bmatrix},
\qquad
D=0.
\]

Il compito consiste nel determinare:
\begin{enumerate}
    \item gli stati raggiungibili dall'origine in \(1,2,3,4\) passi;
    \item gli stati non distinguibili dall'origine in \(1,2,3,4\) passi, cioè il sottospazio di inosservabilità progressiva;
    \item una forma di Kalman del sistema;
    \item gli spazi nei quali evolve l'evoluzione libera per certe condizioni iniziali assegnate.
\end{enumerate}

\subsection{Stati raggiungibili dall'origine in \(1,2,3,4\) passi}

Nel tempo discreto, gli stati raggiungibili in \(t\) passi dall'origine sono generati dalle colonne della matrice

\[
R_t=
\begin{bmatrix}
B & AB & A^2B & \cdots & A^{t-1}B
\end{bmatrix}.
\]

\paragraph{Passo 1.}
Si ha

\[
R_1 = \operatorname{Im}(B)
=
\operatorname{span}\left\{
\begin{bmatrix}
1\\1\\0\\0
\end{bmatrix}
\right\}.
\]

\paragraph{Passo 2.}
Calcoliamo

\[
AB=
\begin{bmatrix}
1\\1\\1\\1
\end{bmatrix}.
\]

Quindi

\[
R_2=
\operatorname{Im}
\begin{bmatrix}
1 & 1\\
1 & 1\\
0 & 1\\
0 & 1
\end{bmatrix}
=
\operatorname{span}\left\{
\begin{bmatrix}
1\\1\\0\\0
\end{bmatrix},
\begin{bmatrix}
0\\0\\1\\1
\end{bmatrix}
\right\}.
\]

Infatti
\[
\begin{bmatrix}
1\\1\\1\\1
\end{bmatrix}
=
\begin{bmatrix}
1\\1\\0\\0
\end{bmatrix}
+
\begin{bmatrix}
0\\0\\1\\1
\end{bmatrix}.
\]

\paragraph{Passo 3.}
Calcoliamo

\[
A^2B=
A(AB)=
A
\begin{bmatrix}
1\\1\\1\\1
\end{bmatrix}
=
\begin{bmatrix}
2\\2\\2\\2
\end{bmatrix}.
\]

Questa colonna è combinazione lineare delle precedenti, quindi

\[
R_3 = R_2.
\]

\paragraph{Passo 4.}
Di conseguenza anche

\[
R_4 = R_3 = R_2.
\]

\paragraph{Conclusione sul sottospazio di raggiungibilità.}
Lo spazio di raggiungibilità è

\[
\mathcal R = R_2 = R_3 = R_4
=
\operatorname{span}\left\{
\begin{bmatrix}
1\\1\\0\\0
\end{bmatrix},
\begin{bmatrix}
0\\0\\1\\1
\end{bmatrix}
\right\},
\]
e quindi
\[
\dim \mathcal R = 2.
\]

Il sistema \emph{non} è completamente raggiungibile, perché lo spazio di stato è \(\mathbb R^4\).

\bigskip

\subsection{Stati non osservabili in \(1,2,3,4\) passi}

Nel tempo discreto, la matrice di osservabilità in \(t\) passi è

\[
O_t=
\begin{bmatrix}
C\\
CA\\
CA^2\\
\vdots\\
CA^{t-1}
\end{bmatrix},
\]
e il sottospazio di inosservabilità in \(t\) passi è

\[
\bar O_t = \ker(O_t).
\]

\paragraph{Passo 1.}
Poiché

\[
C=
\begin{bmatrix}
0 & 1 & 0 & 0
\end{bmatrix},
\]
si ha

\[
\bar O_1
=
\ker(C)
=
\left\{
x\in\mathbb R^4 : x_2=0
\right\}
=
\operatorname{span}\left\{
\begin{bmatrix}
1\\0\\0\\0
\end{bmatrix},
\begin{bmatrix}
0\\0\\1\\0
\end{bmatrix},
\begin{bmatrix}
0\\0\\0\\1
\end{bmatrix}
\right\}.
\]

\paragraph{Passo 2.}
Calcoliamo

\[
CA=
\begin{bmatrix}
0 & 1 & 0 & 1
\end{bmatrix}.
\]

Quindi

\[
O_2=
\begin{bmatrix}
0 & 1 & 0 & 0\\
0 & 1 & 0 & 1
\end{bmatrix}.
\]

Le condizioni \(O_2x=0\) diventano

\[
\begin{cases}
x_2=0,\\
x_2+x_4=0.
\end{cases}
\]

Dunque
\[
x_2=0,\qquad x_4=0,
\]
e quindi

\[
\bar O_2
=
\operatorname{span}\left\{
\begin{bmatrix}
1\\0\\0\\0
\end{bmatrix},
\begin{bmatrix}
0\\0\\1\\0
\end{bmatrix}
\right\}.
\]

\paragraph{Passo 3.}
Calcoliamo

\[
CA^2=
\begin{bmatrix}
1 & 1 & 1 & 1
\end{bmatrix}.
\]

Allora

\[
O_3=
\begin{bmatrix}
0 & 1 & 0 & 0\\
0 & 1 & 0 & 1\\
1 & 1 & 1 & 1
\end{bmatrix}.
\]

Le equazioni \(O_3x=0\) sono

\[
\begin{cases}
x_2=0,\\
x_2+x_4=0,\\
x_1+x_2+x_3+x_4=0.
\end{cases}
\]

Dalle prime due:
\[
x_2=0,\qquad x_4=0.
\]
Sostituendo nella terza:
\[
x_1+x_3=0.
\]

Quindi

\[
\bar O_3
=
\operatorname{span}\left\{
\begin{bmatrix}
1\\0\\-1\\0
\end{bmatrix}
\right\}.
\]

\paragraph{Passo 4.}
Poiché
\[
CA^3=
\begin{bmatrix}
2 & 2 & 2 & 2
\end{bmatrix}
\]
è multiplo di \(CA^2\), non si ottiene alcuna informazione nuova. Dunque

\[
\bar O_4 = \bar O_3.
\]

\paragraph{Conclusione sul sottospazio di inosservabilità.}
Si ha

\[
\bar{\mathcal O}
=
\bar O_3
=
\operatorname{span}\left\{
\begin{bmatrix}
1\\0\\-1\\0
\end{bmatrix}
\right\},
\qquad
\dim \bar{\mathcal O}=1.
\]

\bigskip

\subsection{Forma di Kalman del sistema}

Abbiamo trovato:

\[
\mathcal R
=
\operatorname{span}\left\{
\begin{bmatrix}
1\\1\\0\\0
\end{bmatrix},
\begin{bmatrix}
0\\0\\1\\1
\end{bmatrix}
\right\},
\qquad
\bar{\mathcal O}
=
\operatorname{span}\left\{
\begin{bmatrix}
1\\0\\-1\\0
\end{bmatrix}
\right\}.
\]

Per costruire la forma di Kalman bisogna prima controllare l'intersezione

\[
\mathcal R\cap\bar{\mathcal O}.
\]

Cerchiamo dunque se
\[
\alpha
\begin{bmatrix}
1\\1\\0\\0
\end{bmatrix}
+
\beta
\begin{bmatrix}
0\\0\\1\\1
\end{bmatrix}
=
\gamma
\begin{bmatrix}
1\\0\\-1\\0
\end{bmatrix}.
\]

Confrontando le componenti:

\[
\begin{cases}
\alpha=\gamma,\\
\alpha=0,\\
\beta=-\gamma,\\
\beta=0.
\end{cases}
\]

Ne segue
\[
\alpha=\beta=\gamma=0,
\]
quindi

\[
\mathcal R\cap\bar{\mathcal O}=\{0\}.
\]

Questo vuol dire che \emph{non esiste una parte del sistema che sia contemporaneamente raggiungibile e non osservabile}.

\medskip

Possiamo allora scegliere una base adattata nel modo seguente:
\[
T_{RO}=T_R
=
\begin{bmatrix}
1 & 0\\
1 & 0\\
0 & 1\\
0 & 1
\end{bmatrix},
\qquad
T_{\bar R\bar O}=T_{\bar O}
=
\begin{bmatrix}
1\\
0\\
-1\\
0
\end{bmatrix},
\qquad
T_{\bar R O}=
\begin{bmatrix}
1\\0\\0\\0
\end{bmatrix}.
\]

Una possibile matrice di cambiamento di base è dunque

\[
T=
\begin{bmatrix}
1 & 0 & 1 & 1\\
1 & 0 & 0 & 0\\
0 & 1 & -1 & 0\\
0 & 1 & 0 & 0
\end{bmatrix}.
\]

In queste coordinate il sistema assume una forma di Kalman nella quale compaiono:
\begin{itemize}
    \item un blocco raggiungibile e osservabile di dimensione \(2\);
    \item un blocco non raggiungibile ma osservabile di dimensione \(1\);
    \item un blocco non raggiungibile e non osservabile di dimensione \(1\).
\end{itemize}

Poiché l'intersezione \(\mathcal R\cap\bar{\mathcal O}\) è nulla, il blocco ``raggiungibile e non osservabile'' non compare.

\bigskip

\subsection{Sottospazi di evoluzione libera per alcune condizioni iniziali}

Per un sistema a tempo discreto l'evoluzione libera è

\[
x(k)=A^k x_0.
\]

Il sottospazio generato dall'evoluzione libera a partire da \(x_0\) è il sottospazio ciclico

\[
S(x_0)=\operatorname{span}\{x_0,Ax_0,A^2x_0,\dots\}.
\]

Consideriamo tre condizioni iniziali.

\paragraph{Caso 1:}
\[
x_{0a}=
\begin{bmatrix}
-1\\0\\1\\0
\end{bmatrix}.
\]

Si calcola

\[
Ax_{0a}=
\begin{bmatrix}
1\\0\\-1\\0
\end{bmatrix}
=
-x_{0a}.
\]

Quindi \(x_{0a}\) è autovettore associato all'autovalore \(-1\), e pertanto

\[
S(x_{0a})=\operatorname{span}\{x_{0a}\}.
\]

\paragraph{Caso 2:}
\[
x_{0b}=
\begin{bmatrix}
1\\1\\-1\\-1
\end{bmatrix}.
\]

Si verifica che

\[
Ax_{0b}=0.
\]

Quindi per \(k\ge 1\) l'evoluzione si annulla, e il sottospazio generato è

\[
S(x_{0b})=\operatorname{span}\{x_{0b}\}.
\]

\paragraph{Caso 3:}
\[
x_{0c}=
\begin{bmatrix}
1\\-1\\0\\0
\end{bmatrix}.
\]

Si ottiene

\[
Ax_{0c}=
\begin{bmatrix}
-1\\-1\\1\\1
\end{bmatrix},
\qquad
A^2x_{0c}=0.
\]

Quindi

\[
S(x_{0c})=\operatorname{span}\{x_{0c},Ax_{0c}\}.
\]

Questo terzo caso è importante perché mostra una dinamica \emph{nilpotente}: dopo due passi l'evoluzione libera si annulla.

\bigskip

\section{Esercizio 2: modi del sistema, forma di Jordan e criteri PBH}

Consideriamo ora il sistema

\[
A=
\begin{bmatrix}
-1 & 0 & 1\\
-1 & 0 & 0\\
0 & 0 & -1
\end{bmatrix},
\qquad
B=
\begin{bmatrix}
1\\
0\\
-1
\end{bmatrix},
\qquad
C=
\begin{bmatrix}
-2 & 2 & 0
\end{bmatrix}.
\]

L'esercizio chiede di:
\begin{enumerate}
    \item determinare i modi del sistema;
    \item costruire una base di Jordan;
    \item stabilire quali autovalori sono interni al sottospazio di raggiungibilità e quali sono interni al sottospazio di inosservabilità.
\end{enumerate}

\subsection{Autovalori e struttura di Jordan}

Poiché \(A\) è triangolare a blocchi, i suoi autovalori si leggono facilmente:

\[
\sigma(A)=\{0,-1,-1\}.
\]

Quindi:
\begin{itemize}
    \item l'autovalore \(0\) ha molteplicità algebrica \(1\);
    \item l'autovalore \(-1\) ha molteplicità algebrica \(2\).
\end{itemize}

Studiamo la geometricità di \(-1\). Si ha

\[
A+I=
\begin{bmatrix}
0 & 0 & 1\\
-1 & 1 & 0\\
0 & 0 & 0
\end{bmatrix}.
\]

Il rango è \(2\), quindi

\[
\dim\ker(A+I)=1.
\]

Dunque l'autovalore \(-1\) ha \emph{un solo autovettore} e quindi è associato a un unico blocco di Jordan di ordine \(2\).

Un autovettore è

\[
q_1=
\begin{bmatrix}
1\\1\\0
\end{bmatrix},
\qquad
(A+I)q_1=0.
\]

Serve poi un autovettore generalizzato \(q_2\) tale che

\[
(A+I)q_2=q_1.
\]

Una possibile scelta è

\[
q_2=
\begin{bmatrix}
0\\1\\1
\end{bmatrix},
\]
infatti

\[
(A+I)
\begin{bmatrix}
0\\1\\1
\end{bmatrix}
=
\begin{bmatrix}
1\\1\\0
\end{bmatrix}
=q_1.
\]

Per l'autovalore \(0\) possiamo prendere un autovettore \(q_3\) soluzione di
\[
Aq_3=0.
\]
Ad esempio:

\[
q_3=
\begin{bmatrix}
1\\
1\\
1
\end{bmatrix}.
\]

Quindi una possibile base di Jordan è

\[
T=
\begin{bmatrix}
1 & 0 & 1\\
1 & 1 & 1\\
0 & 1 & 1
\end{bmatrix}.
\]

\bigskip

\subsection{Autovalori interni al sottospazio di raggiungibilità}

Usiamo il criterio PBH per la raggiungibilità.

La coppia \((A,B)\) è raggiungibile su un autovalore \(\lambda\) se

\[
\operatorname{rank}\begin{bmatrix}
A-\lambda I & B
\end{bmatrix}=n.
\]

\paragraph{Caso \(\lambda=0\).}
Si ha

\[
\begin{bmatrix}
A & B
\end{bmatrix}
=
\begin{bmatrix}
-1 & 0 & 1 & 1\\
-1 & 0 & 0 & 0\\
0 & 0 & -1 & -1
\end{bmatrix}.
\]

Il rango è \(2<3\), quindi l'autovalore \(0\) \emph{non è raggiungibile}.

\paragraph{Caso \(\lambda=-1\).}
Si ha

\[
\begin{bmatrix}
A+I & B
\end{bmatrix}
=
\begin{bmatrix}
0 & 0 & 1 & 1\\
-1 & 1 & 0 & 0\\
0 & 0 & 0 & -1
\end{bmatrix},
\]
che ha rango \(3\). Quindi l'autovalore \(-1\) è raggiungibile.

\paragraph{Conclusione.}
L'unico autovalore \emph{interno al sottospazio di raggiungibilità} è

\[
\boxed{\lambda=-1.}
\]

L'autovalore \(0\) è esterno a \(\mathcal R\).

\bigskip

\subsection{Autovalori interni al sottospazio di inosservabilità}

Usiamo ora il criterio PBH per l'osservabilità.

La coppia \((A,C)\) è osservabile su \(\lambda\) se

\[
\operatorname{rank}
\begin{bmatrix}
A-\lambda I\\
C
\end{bmatrix}=n.
\]

\paragraph{Caso \(\lambda=0\).}
Si ottiene

\[
\begin{bmatrix}
A\\
C
\end{bmatrix}
=
\begin{bmatrix}
-1 & 0 & 1\\
-1 & 0 & 0\\
0 & 0 & -1\\
-2 & 2 & 0
\end{bmatrix},
\]
che ha rango \(3\). Quindi \(0\) è osservabile.

\paragraph{Caso \(\lambda=-1\).}
Si ottiene

\[
\begin{bmatrix}
A+I\\
C
\end{bmatrix}
=
\begin{bmatrix}
0 & 0 & 1\\
-1 & 1 & 0\\
0 & 0 & 0\\
-2 & 2 & 0
\end{bmatrix}.
\]

La quarta riga è multiplo della seconda, quindi il rango è \(2<3\). Dunque \(-1\) \emph{non è osservabile}.

\paragraph{Conclusione.}
L'autovalore \(-1\) è interno al sottospazio di inosservabilità:

\[
\boxed{\lambda=-1 \text{ è inosservabile}.}
\]

In sintesi:
\begin{itemize}
    \item \(0\) è osservabile ma non raggiungibile;
    \item \(-1\) è raggiungibile ma non osservabile.
\end{itemize}

Questo è un esempio molto istruttivo, perché mostra che un autovalore può essere presente nel sistema ma non comparire nella parte ingresso--uscita a seconda di come si combinano raggiungibilità e osservabilità.

\bigskip

\section{Esercizio 3: equazione differenziale, realizzazione e modi visibili in uscita}

Consideriamo l'equazione differenziale ingresso--uscita

\[
y^{(4)}(t)+y^{(3)}(t)-\ddot y(t)-\dot y(t)=\dot u(t)-u(t).
\]

\subsection{Funzione di trasferimento}

Applicando la trasformata di Laplace in condizioni iniziali nulle otteniamo

\[
\left(s^4+s^3-s^2-s\right)Y(s)=(s-1)U(s),
\]
da cui

\[
G(s)=\frac{Y(s)}{U(s)}=
\frac{s-1}{s^4+s^3-s^2-s}.
\]

Fattorizzando il denominatore:

\[
s^4+s^3-s^2-s
=
s(s^3+s^2-s-1)
=
s(s+1)^2(s-1).
\]

Quindi

\[
G(s)=\frac{s-1}{s(s+1)^2(s-1)}.
\]

Dal punto di vista della funzione di trasferimento minima si può cancellare il fattore \((s-1)\), ottenendo

\[
G_{\min}(s)=\frac{1}{s(s+1)^2}.
\]

\paragraph{Interpretazione.}
Questo significa che il modo \(\lambda=1\):
\begin{itemize}
    \item è presente in una realizzazione non minima,
    \item ma non compare nella dinamica ingresso--uscita,
    \item quindi è o non raggiungibile, o non osservabile, o entrambe le cose.
\end{itemize}

Nel caso della forma canonica di controllo che costruiremo tra poco, la realizzazione è completamente raggiungibile per costruzione, quindi il modo \(\lambda=1\) sarà necessariamente \emph{non osservabile}.

\bigskip

\subsection{Una realizzazione in forma canonica di controllo}

Poiché il denominatore \emph{non ridotto} è

\[
d(s)=s^4+s^3-s^2-s,
\]
i coefficienti sono

\[
a_0=0,\qquad a_1=-1,\qquad a_2=-1,\qquad a_3=1.
\]

Una forma canonica di controllo è

\[
A_c=
\begin{bmatrix}
0 & 1 & 0 & 0\\
0 & 0 & 1 & 0\\
0 & 0 & 0 & 1\\
0 & 1 & 1 & -1
\end{bmatrix},
\qquad
B_c=
\begin{bmatrix}
0\\0\\0\\1
\end{bmatrix}.
\]

Poiché il numeratore è
\[
n(s)=s-1=-1+s+0s^2+0s^3,
\]
si può scegliere

\[
C_c=
\begin{bmatrix}
-1 & 1 & 0 & 0
\end{bmatrix},
\qquad
D=0.
\]

\paragraph{Osservazione.}
Questa realizzazione è:
\begin{itemize}
    \item completamente raggiungibile, per costruzione;
    \item non completamente osservabile, perché nella funzione di trasferimento è presente una cancellazione polo-zero.
\end{itemize}

Gli autovalori della matrice \(A_c\) sono

\[
0,\,-1,\,-1,\,1.
\]

Dopo la cancellazione nella FDT, resta solo

\[
0,\,-1,\,-1.
\]

Quindi il modo \(\lambda=1\) è un modo interno non osservabile.

\bigskip

\subsection{Quali andamenti possono comparire nell'uscita forzata?}

Dalla funzione di trasferimento minima

\[
G_{\min}(s)=\frac{1}{s(s+1)^2}
\]
deduciamo che i modi ingresso--uscita sono associati ai poli:
\[
0,\quad -1 \text{ con molteplicità } 2.
\]

Quindi gli andamenti fondamentali che possono comparire nell'uscita forzata sono

\[
1,\qquad e^{-t},\qquad t e^{-t}.
\]

Questo perché:
\begin{itemize}
    \item il polo in \(0\) genera il modo costante \(1\);
    \item il polo doppio in \(-1\) genera i modi \(e^{-t}\) e \(t e^{-t}\).
\end{itemize}

\paragraph{Conclusioni sui tre casi tipici.}
\begin{enumerate}
    \item \(\sin t\) e \(1\): \emph{non} può comparire \(\sin t\), perché richiederebbe autovalori complessi puri \(\pm j\), che qui non ci sono. Il modo costante \(1\), invece, è possibile.
    \item \(t e^{-t}\) e \(1\): entrambi sono possibili, perché corrispondono ai poli \(0\) e \(-1\) doppio.
    \item \(e^{t}\) ed \(e^{-t}\): il modo \(e^t\) corrisponde all'autovalore \(1\), ma quel polo è cancellato nella FDT, quindi non è visibile in uscita. Il modo \(e^{-t}\) è invece possibile.
\end{enumerate}

\medskip

La regola generale da ricordare è questa:

\begin{quote}
\emph{nell'uscita forzata compaiono solo i modi associati agli autovalori che sono contemporaneamente raggiungibili e osservabili.}
\end{quote}

\bigskip

\section{Realizzazione dei sistemi e forma minima}

Questa lezione si chiude con una domanda naturale:

\begin{quote}
\emph{dati numeratore e denominatore di una funzione di trasferimento, quale realizzazione in forma di stato conviene scegliere?}
\end{quote}

Sia

\[
G(s)=\frac{n(s)}{d(s)},
\qquad
\deg d(s)=n.
\]

Una \emph{realizzazione minima} di \(G(s)\) è una realizzazione in forma di stato:
\begin{itemize}
    \item con il minimo numero di stati possibile;
    \item completamente raggiungibile;
    \item completamente osservabile.
\end{itemize}

La decomposizione di Kalman spiega esattamente come si ottiene la forma minima: basta eliminare tutti i modi che non sono contemporaneamente raggiungibili e osservabili.

Se il sistema in forma di Kalman è scritto come

\[
\tilde A=
\begin{bmatrix}
A_{RO} & * & * & *\\
0 & A_{R\bar O} & * & *\\
0 & 0 & A_{\bar R O} & *\\
0 & 0 & 0 & A_{\bar R\bar O}
\end{bmatrix},
\qquad
\tilde B=
\begin{bmatrix}
B_{RO}\\
B_{R\bar O}\\
0\\
0
\end{bmatrix},
\qquad
\tilde C=
\begin{bmatrix}
C_{RO} & 0 & C_{\bar R O} & 0
\end{bmatrix},
\]
allora la funzione di trasferimento dipende solo dal blocco
\[
(A_{RO},B_{RO},C_{RO}).
\]

Più precisamente,

\[
G(s)=C_{RO}(sI-A_{RO})^{-1}B_{RO}+D.
\]

Questa è la \emph{forma minima} ottenuta eliminando i modi nascosti dall'ingresso o dall'uscita.

\bigskip

\section{Messaggio finale della lezione}

Questa lezione, anche se fatta in gran parte di esercizi, contiene un messaggio teorico molto importante:

\begin{enumerate}
    \item lo stato completo di un modello può contenere più modi di quelli che la funzione di trasferimento mostra;
    \item la funzione di trasferimento vede solo i modi contemporaneamente raggiungibili e osservabili;
    \item la decomposizione di Kalman è lo strumento che separa sistematicamente:
    \begin{itemize}
        \item parte raggiungibile e osservabile,
        \item parte raggiungibile ma non osservabile,
        \item parte non raggiungibile ma osservabile,
        \item parte non raggiungibile e non osservabile;
    \end{itemize}
    \item la forma minima si ottiene mantenendo solo il sottosistema raggiungibile e osservabile.
\end{enumerate}

\bigskip

\section*{Promemoria operativo per gli esercizi}

Quando in un esercizio ti chiedono di studiare il sistema, la scaletta più efficiente è spesso questa:

\begin{enumerate}
    \item calcolare \(\mathcal R\) oppure \(R_t\) se il sistema è a tempo discreto;
    \item calcolare \(\bar{\mathcal O}\) oppure \(\bar O_t\);
    \item verificare se ci sono intersezioni \(\mathcal R\cap \bar{\mathcal O}\);
    \item usare PBH per associare i singoli autovalori ai sottospazi;
    \item ricostruire la forma di Kalman;
    \item leggere da lì quali modi compaiono nella funzione di trasferimento;
    \item ridurre il sistema alla sua forma minima.
\end{enumerate}